\documentclass[12pt,a4paper]{article}
\usepackage[utf8]{inputenc}
\usepackage[T1]{fontenc}
\usepackage[brazil]{babel}
\usepackage{indentfirst}
\usepackage[margin=2.5cm]{geometry}

\title{Resenha Cr\'itica: Artigo 4 ---\\
Design by Contract}
\author{Gustavo Pessoa Firmino Duarte}
\date{22 de Mar\c{c}o de 2026}

\begin{document}
\maketitle

\section{Introdu\c{c}\~ao}

Bertrand Meyer, criador da linguagem Eiffel, apresenta no artigo ``Design by
Contract'' uma das contribui\c{c}\~oes mais elegantes \`a engenharia de software:
a ideia de que a rela\c{c}\~ao entre um m\'odulo de software e seus clientes pode
--- e deve --- ser formalizada como um contrato jur\'idico. Assim como em um
contrato comercial existem obriga\c{c}\~oes e direitos para ambas as partes, na
programa\c{c}\~ao orientada a objetos cada m\'etodo pode especificar exatamente o
que exige de quem o chama (pr\'e-condi\c{c}\~ao), o que garante ao chamador ap\'os
sua execu\c{c}\~ao (p\'os-condi\c{c}\~ao) e as propriedades que uma classe mant\'em
consistentes em todos os momentos (invariante de classe). Essa abordagem, conhecida
como \textit{Design by Contract} (DbC), transforma a especifica\c{c}\~ao de
software de um exerc\'icio informal em uma ferramenta precisa e verific\'avel.

\section{Contratos, Pr\'e-condi\c{c}\~oes e Invariantes}

O n\'ucleo do DbC repousa em tr\^es elementos. A \textbf{pr\'e-condi\c{c}\~ao}
define as responsabilidades do cliente: se o chamador n\~ao satisfaz a
pr\'e-condi\c{c}\~ao, o comportamento do m\'etodo \'e indefinido e a culpa \'e
do chamador. A \textbf{p\'os-condi\c{c}\~ao} define as responsabilidades do
fornecedor: se a pr\'e-condi\c{c}\~ao foi atendida e o m\'etodo retorna, a
p\'os-condi\c{c}\~ao precisa ser verdadeira. A \textbf{invariante de classe}
especifica propriedades que devem valer ap\'os qualquer cria\c{c}\~ao e ap\'os
qualquer chamada de m\'etodo p\'ublico.

Uma consequ\^encia important\'issima dessa vis\~ao \'e a cr\'itica \`a
\textit{programa\c{c}\~ao defensiva}: Meyer argumenta que verificar a mesma
condi\c{c}\~ao tanto no fornecedor quanto no cliente \'e redundante e contr\'ario
ao esp\'irito dos contratos. Se existe um contrato claro, o fornecedor n\~ao
precisa se proteger de chamadas inv\'alidas --- essas chamadas simplesmente
constituem um erro do cliente. Isso resulta em c\'odigo mais limpo, sem camadas
de valida\c{c}\~ao duplicadas.

Meyer tamb\'em aborda o tratamento de exce\c{c}\~oes atrav\'es dos mecanismos
\textit{rescue} e \textit{retry} do Eiffel: uma exce\c{c}\~ao \'e vista como
uma viola\c{c}\~ao contratual, e a cl\'ausula \textit{rescue} tem a responsabilidade
de tentar restaurar o invariante antes de uma nova tentativa ou de propagar a
falha. Por fim, o princ\'ipio de \textit{subcontra\c{c}\~ao} pela heran\c{c}a ---
subclasses podem apenas afrouxar pr\'e-condi\c{c}\~oes e fortalecer
p\'os-condi\c{c}\~oes, respeitando o Princ\'ipio de Substitui\c{c}\~ao de Liskov
--- integra DbC naturalmente \`a hierarquia de heran\c{c}a.

\section{Conex\~ao com a Pr\'atica Profissional}

Minha experi\^encia como estagi\'ario de Sistemas na ArcelorMittal, atuando no
suporte a aplica\c{c}\~oes internacionais, evidenciou a import\^ancia de contratos
claros entre m\'odulos. Em v\'arios sistemas legados que analisei, a aus\^encia de
pr\'e-condi\c{c}\~oes expl\'icitas gerava um padr\~ao t\'ipico: cada servi\c{c}o
validava entradas que j\'a deveriam ter sido validadas pelo chamador, resultando
em pilhas de \texttt{if (value != null \&\& value.isValid())} espalhadas por toda
a base de c\'odigo. Esse \'e exatamente o antipadr\~ao que Meyer critica ao
defender o DbC contra a programa\c{c}\~ao defensiva.

No meu MVP de e-commerce com Node.js e Prisma ORM, adotei uma pr\'atica
parcialmente inspirada no DbC ao definir esquemas de valida\c{c}\~ao na camada de
API (usando Zod) e tratar esses esquemas como pr\'e-condi\c{c}\~oes: se a requisi\c{c}\~ao
n\~ao passa na valida\c{c}\~ao, a camada de servi\c{c}o nunca \'e acionada. A
camada de servi\c{c}o, por sua vez, assume que os dados est\~ao v\'alidos ---
exatamente a divis\~ao de responsabilidades que Meyer propugna. Essa estrat\'egia
reduziu significativamente a quantidade de \texttt{try/catch} e valida\c{c}\~oes
redundantes no c\'odigo de neg\'ocio.

Reconhe\c{c}o, por\'em, que invariantes de classe s\~ao quase inexistentes em
JavaScript/TypeScript por padr\~ao, e ferramentas como TypeScript apenas oferecem
garantias est\'aticas parciais. Linguagens como Eiffel ou Kotlin com contratos
expl\'icitos representam uma vis\~ao ainda n\~ao amplamente adotada no ecossistema
que utilizo.

\section{Conclus\~ao}

O DbC de Meyer oferece uma mudan\c{c}a de paradigma: em vez de escrever c\'odigo
defensivo que tenta se proteger de todos os abusos poss\'iveis, o desenvolvedor
especifica com precis\~ao o que \'e esperado e o que \'e garantido. Essa clareza
beneficia a documenta\c{c}\~ao, os testes e a manuten\c{c}\~ao. Embora poucos
ambientes modernos suportem DbC nativamente como o Eiffel, os princ\'ipios de
pr\'e e p\'os-condi\c{c}\~oes influenciam fortemente bibliotecas de valida\c{c}\~ao,
frameworks de testes por contrato e at\'e especifica\c{c}\~oes OpenAPI. A ess\^encia
do contrato --- obriga\c{c}\~oes e benef\'icios m\'utuos claramente definidos ---
\'e uma das ideias mais duradouras da engenharia de software.

\end{document}
